\subsection{Regularity repair: total-lapse stages and attenuated coupling}
\label{sec:tele-repair}
The two long-run failures motivate separating the continuum gauge from its near-puncture representation. The telegrapher addition already contains $\chi$ in Eq.~\eqref{eq:telegraph-alpha}. Its logarithmic rate, however, contains $\chi/\alpha$. A collapsing lapse can therefore amplify a finite additive flux source in the stored variable. A lower floor prevents underflow but does not prevent a large positive logarithmic stage increment from exponentiating to a large lapse. The floored terminal replay also finds a neighboring conformal lapse of approximately $7.59\times10^{12}$, consistent with this conditioning problem. This evidence motivates the following separately identified candidate; it does not establish a unique cause of every failed control.

First, the option \code{lapse_update=total} advances the total conformal lapse at ordinary nodes during RK4. Each stage is formed in $\alpha$, receives the same physical-lapse floor, and is converted back to logarithmic storage. The accepted compensated update also uses $\alpha$ units and accounts for reconstruction roundoff. The AF coefficient endpoint and scri retain their previous variables. This changes the finite-step update without changing the underlying lapse equation. The existing filter rate is still computed in the stored chart before conversion to the total-lapse rate.

Second, \code{puncture_stencil=positive} excludes the AF coefficient slot at $r=0$ from positive-radius derivative and advection stencils. The endpoint coefficient is still evolved, but is no longer treated as a neighboring bulk sample. Every positive-radius node remains in the evolution, constraint evaluation, and elliptic solve. This is a finite-grid inner closure, not a derived continuum boundary condition or positive-radius excision. Polynomial stencil identities, endpoint isolation, and independent native/Python rate comparisons test its implementation.

Third, the new input \code{lapse_attenuation} accepts $0\leq\epsilon_\alpha\leq1$ in
\begin{align}
 \eta_\alpha&=\frac{(1+\epsilon_\alpha)\alpha}
 {\alpha+\epsilon_\alpha A},\\
 \Delta(D\alpha)&=\eta_\alpha\chi
 \left[\varpi B_r+\left(\varpi_r+\frac{2\varpi}{r}\right)B\right].
 \label{eq:attenuated-lapse}
\end{align}
Equation~\eqref{eq:telegraph-B} is unchanged. The factor multiplies the existing divergence; it is not moved inside the derivative. For positive $\epsilon_\alpha$, the logarithmic-source coefficient becomes $(1+\epsilon_\alpha)\chi/(\alpha+\epsilon_\alpha A)$ and is bounded as $\alpha\to0$ at fixed finite $\chi$. The reference lapse has $\eta_\alpha=1$; zero attenuation recovers the historical gauge. This third change explicitly modifies the gauge equation. It neither changes a Z4c constraint source nor projects a differential constraint.

The frozen coupled speeds used by both the characteristic diagnostic and timestep guard become
\begin{equation}
 c_\pm=-\beta\pm\sqrt{\chi\left(F+
 \eta_\alpha\frac{\kappa_T\varpi^2}{\tau_T}\right)}.
 \label{eq:attenuated-speeds}
\end{equation}
The frozen dispersion polynomial~\eqref{eq:telegraph-dispersion} retains its form with $c_T^2$ replaced by $\eta_\alpha c_T^2$. Attenuation therefore also reduces the additional frozen-wave damping near lapse collapse. Bounded logarithmic forcing is its purpose; stronger damping everywhere is not claimed. Spatially varying coefficients, the elliptic shift, and nonlinear sources remain outside this frozen argument.

A frozen-symbol check on the saved $t=151.5$ geometry illustrates the distinction between attenuating the coupling and removing its damping. At the first positive node, $r=1.79\times10^{-9}$, the multiplier is $\eta_\alpha=0.00460$, while the high-frequency decay rate from Eq.~\eqref{eq:telegraph-dispersion} is $24.9950$, compared with $25.0000$ when $\eta_\alpha$ is set to unity on the same geometry. The modified $c_+$ is $-2.35\times10^{-9}$; the unattenuated expression on that geometry gives $+1.48\times10^{-9}$. Both modified families point inward at the first seven nodes, but farther-out gauge speeds are positive. This is a local coefficient comparison, not a separately evolved control or a causal-isolation argument for the elliptic system.


\subsection{Failed controls and completed duration evidence}
\label{sec:tele-repair-evidence}
Neither the total-lapse update nor the inner closure alone establishes stability. Table~\ref{tab:repair-controls} records fresh failed controls. A separate blend to derivatives of reconstructed total fields inside $r=10^{-3}$ fails early and is disabled in the current candidate. The two attenuated $N=257$ controls already have approximately seven percent relative horizon-mass error and exterior Hamiltonian RMS of order $0.75$ by $t=31.25$. Their later termination times are not improvements in scientific acceptance.

\begin{table*}[tbp]
\caption{Fresh unsuccessful repair controls, all initialized from the wormhole data. ``Total'' denotes total-lapse RK stages and ``positive'' denotes positive-radius-only bulk stencils. Direct-total derivatives are an additional inner collar experiment. Each failed trajectory is retained under its original input and source binding.}
\label{tab:repair-controls}
\begin{ruledtabular}
\begin{tabular}{p{9.2cm}rrl}
Variant & $N$ & Stop time & Guard\\
\hline
Total &257&17.2758&Speed\\
Total &513&42.6005&Speed\\
Total and positive &257&32.9070&Speed\\
Total, positive, direct-total derivatives &257&4.7414&Positivity\\
Same &513&5.6743&Positivity\\
Total, positive, $\epsilon_\alpha=10^{-4}$ &257&57.1580&Speed\\
Total, positive, $\epsilon_\alpha=10^{-3}$ &257&73.5424&Speed
\end{tabular}
\end{ruledtabular}
\end{table*}

The present duration candidate combines total-lapse RK, positive-radius stencils, $\epsilon_\alpha=10^{-4}$, and physical floor $10^{-10}$. The direct-total derivative collar is disabled. All other selected mapped-grid, telegrapher, timestep, and KO parameters are retained. No shift cap, upper lapse cap, constraint projection, or relaxed timestep guard is introduced. Nine targeted implementation test groups pass, including independent total-lapse RK stages, exact compensated restart, endpoint isolation, and attenuated native/Python rates and characteristic speeds. These tests do not replace the duration experiment.

Fresh $N=513$ and $1025$ trials target $t=500$ in the compactification units defined in Sec.~\ref{sec:formulation}. The $N=513$ trial is exploratory; the predeclared principal $N=1025$ budget requires relative horizon-mass error below $10^{-5}$. Its full saved-trajectory exterior RMS limits for $(\HH,\MM,\Th,Z_r)$ are $(5\times10^{-5},10^{-4},5\times10^{-7},5\times10^{-6})$; for $t\geq10$ the first two tighten to $10^{-6}$ and $10^{-5}$. The corresponding clock, nullness, and all-stage scri-pole limits are $10^{-12}$, $10^{-10}$, and $10^{-5}$. The interior is excluded and the radial weights are renormalized as in Eq.~\eqref{eq:exterior-norm}. Duration and accuracy are separate requirements.

The principal $N=1025$ run finishes normally at $t=500$ after 10\,240\,000 accepted steps. All 10\,001 saved states are checked for frozen source bindings, finite arrays, reconstructed full-domain norms, and independently computed apparent horizons; the final checkpoint agrees exactly with the last profile. The unchanged predeclared budget passes over the complete saved trajectory. Final relative horizon-mass error is $1.15\times10^{-6}$, with an initial-transient maximum of $2.64\times10^{-6}$. Late exterior RMS maxima are $(1.79\times10^{-7},2.44\times10^{-6},5.00\times10^{-8},5.21\times10^{-7})$. The maximum mapped Courant number is $0.2000$, $\Delta t/\tau_T=0.00244$, and the maximum recorded scri-pole residue is $2.83\times10^{-6}$. Thus the candidate completes the requested duration without relaxing its guards or excluding any positive-radius evolution node.

The exploratory $N=513$ run also completes $t=500$ and validates all 10\,001 saved states, but its final and maximum relative horizon-mass error is $1.09\times10^{-3}$ ($0.109\%$). Its maximum scri-pole residue is $3.35\times10^{-5}$. Table~\ref{tab:repair-duration} reports the final exterior errors at both resolutions. The coarse run is duration-survival evidence, not the principal accuracy certificate; these two runs do not constitute a convergence study of the modified gauge and inner closure.

\begin{table*}[tbp]
\caption{Completed attenuation trials at $t=500$. The horizon-mass column is $|M_{\rm AH}-M|/M$; the remaining columns are exterior RMS norms with the normalization of Eq.~\eqref{eq:exterior-norm}. The $N=1025$ budget is evaluated over the full saved trajectory, not only these final values. The $N=513$ row is an exploratory resolution control.}
\label{tab:repair-duration}
\begin{ruledtabular}
\begin{tabular}{rrrrrr}
$N$ & Horizon-mass error & $\HH$ & $\MM$ & $\Th$ & $Z_r$\\
\hline
513 & $1.09\times10^{-3}$ & $9.64\times10^{-6}$ & $1.27\times10^{-4}$ & $2.61\times10^{-6}$ & $2.71\times10^{-5}$\\
1025 & $1.15\times10^{-6}$ & $1.79\times10^{-7}$ & $2.44\times10^{-6}$ & $5.00\times10^{-8}$ & $5.21\times10^{-7}$
\end{tabular}
\end{ruledtabular}
\end{table*}

Figure~\ref{fig:repair-duration} retains the small inner lapse ripples. Over $0<r\leq10^{-4}$, the maximum fractional physical-lapse change between $t=25$ and $500$ is $3.08\times10^{-4}$. The minimum at $t=500$ is $4.62\times10^{-7}$, above the imposed floor. Five negative neighboring differences remain, with a largest local fractional drop of $0.250$. These observations support settling of this discrete profile, not elimination of its inner oscillation. Interior pointwise errors are still large: the final Hamiltonian maximum is about $393$ at $r=9.52\times10^{-9}$. They remain recorded separately and excluded from the stated exterior budget.

Floor accounting separates initialization from later updates. The initial floor changes 80 positive-radius lapse nodes, all at $r\leq6.29\times10^{-7}$, with maximum conformal-lapse change $5.00\times10^{-11}$. The AF coefficient, scri lapse, and other initial fields are unchanged. The native absolute-correction counter covers subsequent total-lapse RK stages and accepted updates; its maximum is $5.74\times10^{-15}$, with no recorded nonpositive total-lapse clamp. It does not include the initial logarithmic flooring operation. The initial logarithmic correction reaches $11.73$, explaining why logarithmic and absolute correction sizes must be distinguished.

The long-time ray diagnostic also reaches a representation limit. At $t=500$ the 81 inward-family ray coordinates have coalesced at the subnormal binary64 radius $3.72\times10^{-320}$, with minimum stored $\log J\simeq-798$ and zero resolved neighboring gaps. These passive diagnostic rays do not feed back into the field evolution. Their loss of separation is not evidence of a resolved continuum crossing, and it prevents a long-time noncrossing claim from this packet. Finite-duration survival and the exterior accuracy pass therefore do not establish continuum stability or convergence of the regularization.

\begin{figure*}[tbp]
 \centering
 \includegraphics[width=\textwidth]{telegraph_repair_duration}
 \caption{Completed $N=1025$ attenuation candidate through $t=500$. (a) Exterior RMS constraints, including the initial transient. (b) Relative apparent-horizon mass error and its declared limit. (c) Raw physical-lapse profiles near the puncture, including the initial floor and residual late ripples; the $t=25$ and $500$ curves nearly coincide. (d) Conformal lapse across the domain; the dotted line marks the final horizon. No radial smoothing is applied. The full saved trajectory passes its predeclared duration budget. This result does not extend the earlier three-grid convergence study to the modified gauge.}
 \label{fig:repair-duration}
\end{figure*}
